
\documentclass[11pt,a4paper]{article}

\usepackage[margin=1in]{geometry}
\usepackage{amsmath,amssymb,amsthm,mathtools}
\usepackage{hyperref}
\usepackage{booktabs}
\usepackage{microtype}

\hypersetup{
    colorlinks=true,
    linkcolor=blue,
    urlcolor=blue,
    citecolor=blue
}

\title{On the Identity Between the Sum of Cubes and the Square of the Sum}
\author{Dakshesh Nankani}
\date{November 2026}

\begin{document}

\maketitle

\section*{Proposition}

For every positive integer $n$,
\[
\sum_{k=1}^{n} k^3
=
\left(\sum_{k=1}^{n} k\right)^2.
\]

Equivalently,
\[
\sum_{k=1}^{n} k^3
=
\frac{n^2(n+1)^2}{4}.
\]

We establish this identity through two independent approaches:
generating functions, and symmetric polynomials using Newton's
identities.

\section*{Proof I. Generating Functions}

Let
\[
G(x)=\sum_{k=1}^{\infty}k^3x^k
\]
denote the ordinary generating function of the cubes.

As derived in \hyperref[appendix:generating-functions]{Appendix A},
\[
G(x)=\frac{x(1+4x+x^2)}{(1-x)^4}.
\]

Now let
\[
S_n=\sum_{k=1}^{n}k^3
\]
and consider the ordinary generating function of these partial sums:
\[
F(x)=\sum_{n=1}^{\infty}S_nx^n.
\]

By the definition of $S_n$,
\[
F(x)=
\sum_{n=1}^{\infty}
\left(\sum_{k=1}^{n}k^3\right)x^n.
\]

Interchanging the order of summation,
\[
\begin{aligned}
F(x)
&=\sum_{k=1}^{\infty}k^3
\sum_{n=k}^{\infty}x^n\\
&=\sum_{k=1}^{\infty}k^3
\frac{x^k}{1-x}\\
&=\frac{1}{1-x}
\sum_{k=1}^{\infty}k^3x^k.
\end{aligned}
\]

Recognising the inner summation as $G(x)$, we obtain
\[
F(x)=\frac{G(x)}{1-x}.
\]

Consequently,
\[
F(x)=\frac{x(1+4x+x^2)}{(1-x)^5}.
\]

Using the binomial-series expansion,
\[
\frac{1}{(1-x)^5}
=
\sum_{r=0}^{\infty}\binom{r+4}{4}x^r,
\]
we obtain
\[
F(x)=
x(1+4x+x^2)
\sum_{r=0}^{\infty}\binom{r+4}{4}x^r.
\]

Comparing coefficients of $x^n$,
\[
S_n=
\binom{n+3}{4}
+4\binom{n+2}{4}
+\binom{n+1}{4},
\]
where $\binom{m}{4}=0$ for non-negative integers $m<4$.

As shown in \hyperref[appendix:binomial-simplification]{Appendix B}, this expression simplifies to
\[
S_n=\frac{n^2(n+1)^2}{4}.
\]

Therefore,
\[
\boxed{
\sum_{k=1}^{n}k^3
=
\left(\sum_{k=1}^{n}k\right)^2
}.
\]

% \hfill $\square$

\section*{Proof II. Symmetric Polynomials and Newton's Identities}

Let $x_1,x_2,\ldots,x_n$ be indeterminates. Define the
elementary symmetric polynomials by
\[
e_r=
\sum_{1\leq i_1<\cdots<i_r\leq n}
x_{i_1}x_{i_2}\cdots x_{i_r},
\]
and the power sum symmetric polynomials by
\[
p_r=\sum_{i=1}^{n}x_i^r.
\]

The first three Newton identities are
\[
p_1=e_1,
\]
\[
p_2=e_1p_1-2e_2,
\]
\[
p_3=e_1p_2-e_2p_1+3e_3.
\]

Consequently,
\[
p_2=e_1^2-2e_2
\]
and
\[
\boxed{p_3=e_1^3-3e_1e_2+3e_3}.
\]

We now specialise to $x_k=k$, for $1\leq k\leq n$.

The first elementary symmetric polynomial is
\[
e_1=\sum_{k=1}^{n}k
=\frac{n(n+1)}{2}.
\]

The second and third elementary symmetric polynomials,
as derived independently in \hyperref[appendix:symmetric-polynomials]{Appendix C}, are
\[
e_2=
\frac{n(n+1)(n-1)(3n+2)}{24}
\]
and
\[
e_3=
\frac{n^2(n-1)(n-2)(n+1)^2}{48}.
\]

Substituting these expressions into
\[
p_3=e_1^3-3e_1e_2+3e_3,
\]
we obtain
\[
\begin{aligned}
p_3
={}&\frac{n^3(n+1)^3}{8}\\
&-\frac{n^2(n+1)^2(n-1)(3n+2)}{16}\\
&+\frac{n^2(n-1)(n-2)(n+1)^2}{16}.
\end{aligned}
\]

Factoring and simplifying,
\[
\begin{aligned}
p_3
&=\frac{n^2(n+1)^2}{16}
\left[
2n(n+1)
-(n-1)(3n+2)
+(n-1)(n-2)
\right]\\
&=\frac{n^2(n+1)^2}{16}\cdot 4\\
&=\frac{n^2(n+1)^2}{4}.
\end{aligned}
\]

Since $p_3=\sum_{k=1}^{n}k^3$, it follows that
\[
\boxed{
\sum_{k=1}^{n}k^3
=
\left(\sum_{k=1}^{n}k\right)^2
}.
\]

% \hfill $\square$

\appendix

\section{Derivation of the Generating Function for Cubes}
\label{appendix:generating-functions}

Consider the geometric series
\[
\sum_{k=0}^{\infty}x^k=\frac{1}{1-x}.
\]

Define the differential operator
\[
D=x\frac{d}{dx}.
\]

Since
\[
D(x^k)=kx^k,
\]
repeated application of $D$ gives
\[
D^3\left(\sum_{k=0}^{\infty}x^k\right)
=
\sum_{k=1}^{\infty}k^3x^k.
\]

We evaluate the right-hand side successively:
\[
D\left(\frac{1}{1-x}\right)
=
\frac{x}{(1-x)^2},
\]
\[
D^2\left(\frac{1}{1-x}\right)
=
\frac{x(1+x)}{(1-x)^3},
\]
and
\[
D^3\left(\frac{1}{1-x}\right)
=
\frac{x(1+4x+x^2)}{(1-x)^4}.
\]

Therefore, the ordinary generating function for the cubes is
\[
\boxed{
G(x)=\sum_{k=1}^{\infty}k^3x^k
=\frac{x(1+4x+x^2)}{(1-x)^4}
}.
\]

\section{Simplification of the Binomial-Coefficient Expression}
\label{appendix:binomial-simplification}

From Proof I, we obtained
\[
S_n=
\binom{n+3}{4}
+4\binom{n+2}{4}
+\binom{n+1}{4}.
\]

Using the factorial definition of binomial coefficients,
\[
\begin{aligned}
S_n
={}&\frac{(n+3)(n+2)(n+1)n}{24}\\
&+\frac{4(n+2)(n+1)n(n-1)}{24}\\
&+\frac{(n+1)n(n-1)(n-2)}{24}.
\end{aligned}
\]

Factoring out the common terms,
\[
\begin{aligned}
S_n
=\frac{n(n+1)}{24}\big[
&(n+3)(n+2)\\
&+4(n+2)(n-1)\\
&+(n-1)(n-2)
\big].
\end{aligned}
\]

Expanding the expression inside the brackets,
\[
\begin{aligned}
S_n
&=\frac{n(n+1)}{24}
\left[
n^2+5n+6
+4n^2+4n-8
+n^2-3n+2
\right]\\
&=\frac{n(n+1)}{24}(6n^2+6n)\\
&=\frac{n^2(n+1)^2}{4}.
\end{aligned}
\]

Hence,
\[
\boxed{
S_n=\frac{n^2(n+1)^2}{4}
}.
\]

\section{Evaluation of the Second and Third Elementary Symmetric Polynomials}
\label{appendix:symmetric-polynomials}

Let
\[
P_n(t)=\prod_{k=1}^{n}(1+kt)
=\sum_{r=0}^{n}e_r(n)t^r,
\]
where $e_r(n)$ denotes the $r$-th elementary symmetric
polynomial evaluated at $1,2,\ldots,n$.

Observe that
\[
P_n(t)=(1+nt)P_{n-1}(t).
\]

Comparing coefficients of $t^r$ gives the recurrence
\[
e_r(n)=e_r(n-1)+ne_{r-1}(n-1).
\]

\subsection*{Evaluation of $e_2$}

For $r=1$, the recurrence gives
\[
e_1(n)=e_1(n-1)+n.
\]

Since
\[
e_1(n)=\frac{n(n+1)}{2},
\]
the recurrence for $e_2$ becomes
\[
\begin{aligned}
e_2(n)
&=e_2(n-1)+ne_1(n-1)\\
&=e_2(n-1)+\frac{n^2(n-1)}{2}.
\end{aligned}
\]

We claim that
\[
e_2(n)=
\frac{n(n+1)(n-1)(3n+2)}{24}.
\]

To verify this expression, let
\[
E_2(n)=\frac{n(n+1)(n-1)(3n+2)}{24}.
\]

For $n\geq 2$,
\[
\begin{aligned}
E_2(n)-E_2(n-1)
&=\frac{n(n-1)}{24}
\left[
(n+1)(3n+2)-(n-2)(3n-1)
\right]\\
&=\frac{n^2(n-1)}{2}\\
&=ne_1(n-1).
\end{aligned}
\]

Furthermore, $E_2(1)=0=e_2(1)$. Hence, by the recurrence,
\[
\boxed{
e_2(n)=
\frac{n(n+1)(n-1)(3n+2)}{24}
}.
\]

\subsection*{Evaluation of $e_3$}

For $r=3$, the recurrence gives
\[
e_3(n)=e_3(n-1)+ne_2(n-1).
\]

Using the expression for $e_2(n-1)$,
\[
e_2(n-1)
=
\frac{n(n-1)(n-2)(3n-1)}{24},
\]
we obtain
\[
e_3(n)=e_3(n-1)
+\frac{n^2(n-1)(n-2)(3n-1)}{24}.
\]

Since $e_3(2)=0$, we claim that
\[
e_3(n)=
\frac{n^2(n-1)(n-2)(n+1)^2}{48}.
\]

Let
\[
E_3(n)=
\frac{n^2(n-1)(n-2)(n+1)^2}{48}.
\]

We verify this expression using the recurrence. For $n\geq 3$,
\[
\begin{aligned}
E_3(n)-E_3(n-1)
&=\frac{n^2(n-2)}{48}
\left[
(n-1)(n+1)^2-(n-1)^2(n-3)
\right]\\
&=\frac{n^2(n-1)(n-2)}{48}
\left[
(n+1)^2-(n-1)(n-3)
\right]\\
&=\frac{n^2(n-1)(n-2)(3n-1)}{24}\\
&=ne_2(n-1).
\end{aligned}
\]

Furthermore, $E_3(2)=0=e_3(2)$. Hence, by the recurrence,
\[
\boxed{
e_3(n)=
\frac{n^2(n-1)(n-2)(n+1)^2}{48}
}.
\]

\end{document}
